\chapter{KHỞI TẠO QUẦN THỂ THÔNG MINH}
\label{chap:init}

\section{Mục tiêu và yêu cầu}

Thuật toán gốc tạo $N_{pop}$ hoán vị ngẫu nhiên. Với cách cắt tuyến \code{capacity\_time}, một hoán vị
ngẫu nhiên thường đi vòng xa và mở nhiều xe, nên quần thể ban đầu nằm rất xa vùng tốt. Cải tiến
\code{init\_smart} thay bước này bằng ba nguồn cá thể:
\begin{itemize}
  \item một phần quần thể dựng từ thứ tự láng giềng gần nhất (nhãn \code{nn});
  \item một phần dựng từ thứ tự theo khung giờ (nhãn \code{tw});
  \item phần còn lại sinh ngẫu nhiên đúng như thuật toán gốc (nhãn \code{random}).
\end{itemize}
Các yêu cầu cụ thể của cải tiến:
\begin{enumerate}
  \item mọi cá thể là hoán vị hợp lệ của $1,\dots,n$, không thiếu, không trùng khách, không chứa kho;
  \item số cá thể của từng nguồn đúng theo tỉ lệ trong cấu hình, kể cả khi $N_{pop}$ không chia hết;
  \item hai thứ tự NN và TW là tất định, nên cần tạo biến thể để quần thể không chứa nhiều bản giống hệt
        nhau;
  \item chỉ dùng bộ sinh số ngẫu nhiên được truyền vào, cùng seed thì ra cùng quần thể;
  \item khi cờ tắt, \code{nsga2.py} không gọi đến mô-đun này.
\end{enumerate}

Tham số trong \code{config.yaml}:
\begin{Verbatim}
improvements:
  init_smart: false
  init_smart_params:
    ratios: {nn: 0.2, tw: 0.2, random: 0.6}
    perturb_swaps_min: 1
    perturb_swaps_max: 3
    keep_pure_seed: true
    dedup_max_attempts: 50
\end{Verbatim}

\section{Chia số lượng cá thể theo tỉ lệ}
\label{sec:quota}

Với $N_{pop} = 80$ và tỉ lệ $0{,}2/0{,}2/0{,}6$, số lượng là $16/16/48$. Khi $N_{pop}$ không chia hết,
làm tròn từng nguồn riêng lẻ có thể ra tổng khác $N_{pop}$. Mô-đun dùng phương pháp phần dư lớn nhất:
\begin{enumerate}
  \item với mỗi nguồn $k$, lấy phần nguyên $\lfloor N_{pop}\, r_k \rfloor$;
  \item số chỗ còn thiếu bằng $N_{pop}$ trừ tổng các phần nguyên;
  \item phát lần lượt từng chỗ còn thiếu cho nguồn có phần dư $N_{pop} r_k - \lfloor N_{pop} r_k
        \rfloor$ lớn nhất; khi hai phần dư bằng nhau thì ưu tiên theo thứ tự \code{nn}, \code{tw},
        \code{random}.
\end{enumerate}
Tổng luôn đúng bằng $N_{pop}$ và mỗi nguồn lệch tỉ lệ lý tưởng không quá một cá thể.

\begin{table}[H]
\centering
\caption{Ví dụ chia số lượng với tỉ lệ $0{,}2/0{,}2/0{,}6$}
\label{tab:quota}
\begin{tabular}{@{}rccccc@{}}
\toprule
$N_{pop}$ & $N_{pop}\cdot r$ & Phần nguyên & Phần dư & Chỗ thiếu & Kết quả \\
\midrule
80 & 16 / 16 / 48 & 16 / 16 / 48 & 0 / 0 / 0 & 0 & 16 / 16 / 48 \\
25 & 5 / 5 / 15 & 5 / 5 / 15 & 0 / 0 / 0 & 0 & 5 / 5 / 15 \\
7 & 1,4 / 1,4 / 4,2 & 1 / 1 / 4 & 0,4 / 0,4 / 0,2 & 1 & 2 / 1 / 4 \\
3 & 0,6 / 0,6 / 1,8 & 0 / 0 / 1 & 0,6 / 0,6 / 0,8 & 2 & 1 / 0 / 2 \\
\bottomrule
\end{tabular}
\end{table}

Ở dòng $N_{pop} = 7$, \code{nn} và \code{tw} cùng phần dư $0{,}4$; quy tắc phá hoà cho chỗ thừa vào
\code{nn}. Ở dòng $N_{pop} = 3$, hai chỗ thừa vào \code{random} (dư $0{,}8$) rồi \code{nn} (dư $0{,}6$,
thắng \code{tw} khi hoà).

Tham số tỉ lệ được kiểm tra trước khi dùng: đủ ba khoá, mỗi giá trị là số hữu hạn không âm, tổng bằng 1
với sai số $10^{-6}$. Cấu hình sai báo lỗi \code{ValueError} kèm lý do, không âm thầm sửa.

\section{Hai thứ tự gốc tất định}

\subsection{Thứ tự láng giềng gần nhất}

Bắt đầu tại kho (điểm 0). Ở mỗi bước, trong các khách chưa ghé, chọn khách $c$ có $d_{\text{hiện tại},c}$
nhỏ nhất; nếu bằng nhau thì chọn khách có số hiệu nhỏ hơn. Khung thời gian không được xét ở bước này;
việc cắt tuyến và tính giờ đến do bộ giải mã đảm nhiệm.

\subsection{Thứ tự theo khung giờ}

Chọn lần lượt khách chưa ghé có $ETa$ nhỏ nhất. Nếu nhiều khách cùng $ETa$, chọn khách gần điểm vừa chọn
nhất (với khách đầu tiên là gần kho nhất); nếu vẫn bằng nhau thì chọn số hiệu nhỏ hơn. Mốc sắp xếp là
$ETa$ (đầu khung kỳ vọng), không phải $ET$ (đầu khung chấp nhận được): $ET$ bị nới thêm 30 phút và kẹp
về 0, nên nhiều khách có cùng $ET = 0$ dù khung kỳ vọng khác nhau.

\subsection{Ví dụ trên 8 khách đầu của C102}

Bảng~\ref{tab:c102-8} là dữ liệu thật của 8 khách đầu tiên trong C102.

\begin{table}[H]
\centering
\caption{Dữ liệu 8 khách đầu của instance C102}
\label{tab:c102-8}
\begin{tabular}{@{}crrrrr@{}}
\toprule
Khách & $x$ & $y$ & Nhu cầu (kg) & $[ETa, LTb]$ (phút) & $d_{0,i}$ (km) \\
\midrule
0 (kho) & 40 & 50 & 0 & [0, 1236] & 0 \\
1 & 45 & 68 & 10 & [0, 1127] & 18,68 \\
2 & 45 & 70 & 30 & [0, 1125] & 20,62 \\
3 & 42 & 66 & 10 & [0, 1129] & 16,12 \\
4 & 42 & 68 & 10 & [727, 782] & 18,11 \\
5 & 42 & 65 & 10 & [0, 1130] & 15,13 \\
6 & 40 & 69 & 20 & [621, 702] & 19,00 \\
7 & 40 & 66 & 20 & [0, 1130] & 16,00 \\
8 & 38 & 68 & 20 & [255, 324] & 18,11 \\
\bottomrule
\end{tabular}
\end{table}

\paragraph{Thứ tự NN.} Từ kho, khách 5 gần nhất (15,13 km). Từ khách 5, khách 3 cách 1 km. Từ khách 3,
khách 4 và khách 7 cùng cách 2 km, quy tắc phá hoà chọn khách 4. Tiếp theo là 6, 8, 7, 1, 2. Kết quả
$\pi_{NN} = (5, 3, 4, 6, 8, 7, 1, 2)$.

\paragraph{Thứ tự TW.} Năm khách 1, 2, 3, 5, 7 cùng $ETa = 0$. Trong nhóm này, khách 5 gần kho nhất, sau
đó lần lượt chọn khách gần điểm vừa chọn: 3, 7, 1, 2. Tiếp theo là khách 8 ($ETa = 255$), khách 6
($ETa = 621$), khách 4 ($ETa = 727$). Kết quả $\pi_{TW} = (5, 3, 7, 1, 2, 8, 6, 4)$.

Hình~\ref{fig:nn-tw} vẽ hai thứ tự này. Thứ tự NN đi vòng gọn trong cụm khách, còn thứ tự TW quay lại
các khách 8, 6, 4 ở cuối vì khung giờ của chúng mở muộn.

\begin{figure}[H]
\centering
\begin{tikzpicture}
\begin{groupplot}[group style={group size=2 by 1, horizontal sep=1.6cm},
  width=6.6cm, height=6.6cm, xmin=36.5, xmax=46.5, ymin=63.5, ymax=71.5,
  xlabel={$x$}, ylabel={$y$}, tick label style={font=\footnotesize},
  title style={font=\small}]
\nextgroupplot[title={Láng giềng gần nhất}]
\addplot[-{Stealth[length=2mm]}, thick, blue1] coordinates
  {(42,65) (42,66) (42,68) (40,69) (38,68) (40,66) (45,68) (45,70)};
\addplot[only marks, mark=*, mark size=2.2pt, black] coordinates
  {(45,68) (45,70) (42,66) (42,68) (42,65) (40,69) (40,66) (38,68)};
\node[font=\scriptsize, anchor=west] at (axis cs:45.2,68) {1};
\node[font=\scriptsize, anchor=west] at (axis cs:45.2,70) {2};
\node[font=\scriptsize, anchor=west] at (axis cs:42.2,66) {3};
\node[font=\scriptsize, anchor=west] at (axis cs:42.2,68.2) {4};
\node[font=\scriptsize, anchor=west] at (axis cs:42.2,64.8) {5};
\node[font=\scriptsize, anchor=south] at (axis cs:40,69.2) {6};
\node[font=\scriptsize, anchor=north] at (axis cs:40,65.8) {7};
\node[font=\scriptsize, anchor=east] at (axis cs:37.8,68) {8};
\nextgroupplot[title={Theo khung giờ $ETa$}]
\addplot[-{Stealth[length=2mm]}, thick, orange1] coordinates
  {(42,65) (42,66) (40,66) (45,68) (45,70) (38,68) (40,69) (42,68)};
\addplot[only marks, mark=*, mark size=2.2pt, black] coordinates
  {(45,68) (45,70) (42,66) (42,68) (42,65) (40,69) (40,66) (38,68)};
\node[font=\scriptsize, anchor=west] at (axis cs:45.2,68) {1};
\node[font=\scriptsize, anchor=west] at (axis cs:45.2,70) {2};
\node[font=\scriptsize, anchor=west] at (axis cs:42.2,66) {3};
\node[font=\scriptsize, anchor=west] at (axis cs:42.2,68.2) {4};
\node[font=\scriptsize, anchor=west] at (axis cs:42.2,64.8) {5};
\node[font=\scriptsize, anchor=south] at (axis cs:40,69.2) {6};
\node[font=\scriptsize, anchor=north] at (axis cs:40,65.8) {7};
\node[font=\scriptsize, anchor=east] at (axis cs:37.8,68) {8};
\end{groupplot}
\end{tikzpicture}
\caption{Thứ tự ghé của hai heuristic trên 8 khách đầu của C102 (kho ở toạ độ $(40, 50)$, nằm ngoài khung
hình, cả hai thứ tự đều bắt đầu từ khách 5)}
\label{fig:nn-tw}
\end{figure}

Để thấy hai thứ tự gốc khác hoán vị ngẫu nhiên thế nào, bảng~\ref{tab:base-obj} chấm ba hoán vị đầy đủ
của C102-25 bằng bộ giải mã của nhóm.

\begin{table}[H]
\centering
\caption{Mục tiêu của ba hoán vị trên C102-25 (hoán vị ngẫu nhiên lấy với seed 0)}
\label{tab:base-obj}
\begin{tabular}{@{}lrrrr@{}}
\toprule
Hoán vị & $f_1$ (CNY) & $f_2$ (kg) & $f_3$ & Số xe \\
\midrule
Ngẫu nhiên & 4817,5 & 485,2 & 0,4000 & 11 \\
Láng giềng gần nhất & \textbf{2564,4} & \textbf{224,2} & \textbf{0,7034} & 7 \\
Theo khung giờ & 4170,1 & 469,5 & 0,6726 & 10 \\
\bottomrule
\end{tabular}
\end{table}

Thứ tự NN cho chi phí và khí thải chỉ bằng khoảng một nửa hoán vị ngẫu nhiên, thứ tự TW cho mức hài lòng
cao gần bằng NN. Đây là lý do dùng cả hai: mỗi thứ tự gieo một điểm xuất phát tốt ở một vùng khác của
mặt Pareto.

\section{Tạo biến thể bằng hoán đổi}

Hai thứ tự gốc luôn cho cùng một hoán vị, nên nếu sao chép 16 lần thì quần thể có 16 cá thể giống hệt
nhau. Mô-đun giữ một bản gốc (khi \code{keep\_pure\_seed} bật) và tạo các bản còn lại bằng $k$ lần hoán
đổi hai vị trí khác nhau, với $k$ bốc đều trong $[\code{perturb\_swaps\_min},
\code{perturb\_swaps\_max}] = [1, 3]$. Phép hoán đổi giữ nguyên tập khách nên kết quả vẫn là hoán vị hợp
lệ. Ba biến thể đầu tiên của $\pi_{NN}$ trên C102-8 với seed 0 là:
\begin{center}
\begin{tabular}{@{}ll@{}}
Gốc & $(5, 3, 4, 6, 8, 7, 1, 2)$ \\
Biến thể 1 & $(3, 4, 5, 6, 2, 7, 1, 8)$ \\
Biến thể 2 & $(5, 3, 4, 6, 1, 8, 7, 2)$ \\
Biến thể 3 & $(5, 3, 4, 2, 8, 6, 1, 7)$ \\
\end{tabular}
\end{center}
Số lần hoán đổi nhỏ giữ lại phần lớn cấu trúc tốt của thứ tự gốc mà vẫn tạo được khác biệt.

\section{Khử trùng và trường hợp đặc biệt}

Mỗi cá thể mới được so với tập các hoán vị đã có (lưu dưới dạng bộ giá trị để tra cứu trong $O(1)$). Nếu
trùng, mô-đun sinh lại, tối đa \code{dedup\_max\_attempts} lần. Hết số lần thử mà vẫn trùng thì chỗ đó
được thay bằng một hoán vị ngẫu nhiên, gắn nhãn riêng \code{random\_fallback} và được đếm trong
\code{effective\_counts}. Như vậy báo cáo luôn biết chính xác quần thể thực có bao nhiêu cá thể mỗi loại,
thay vì giả định tỉ lệ đã đạt.

Với $n$ rất nhỏ, số hoán vị khác nhau là $n!$ có thể nhỏ hơn $N_{pop}$ (ví dụ $n = 4$ chỉ có 24 hoán
vị). Khi đó không thể có quần thể toàn cá thể khác nhau. Mô-đun vẫn khử trùng cho tới khi đã dùng hết
các hoán vị khả dĩ, sau đó chấp nhận trùng, ghi số cá thể trùng vào \code{n\_duplicates} và một dòng
ghi chú. Hàm không bao giờ lặp vô hạn. Với $n = 1$, phép hoán đổi không thực hiện được và mô-đun trả bản
sao mà không bốc số ngẫu nhiên.

\section{Thuật toán và quy tắc dùng số ngẫu nhiên}

\begin{algorithm}[H]
\caption{Khởi tạo quần thể thông minh}
\label{alg:init}
\KwIn{instance, $N_{pop}$, bộ sinh \textit{rng}, tham số, hàm sinh cá thể ngẫu nhiên của thuật toán gốc \textit{base}}
\KwOut{danh sách $N_{pop}$ hoán vị và nhãn nguồn}
$(q_{nn}, q_{tw}, q_{rd}) \gets$ \textsc{ChiaSoLuong}($N_{pop}$, tỉ lệ)\;
$P \gets \emptyset$; $\textit{seen} \gets \emptyset$\;
\ForEach{$(\textit{nhãn}, \pi_0) \in \{(\texttt{nn}, \pi_{NN}), (\texttt{tw}, \pi_{TW})\}$}{
  \For{$j \gets 1$ \KwTo $q_{\textit{nhãn}}$}{
    \eIf{$j = 1$ \textbf{and} giữ bản gốc \textbf{and} $\pi_0 \notin \textit{seen}$}{thêm $\pi_0$ với nhãn \textit{nhãn}\;}{
      thử tối đa $A$ lần: $\pi \gets$ \textsc{HoanDoi}($\pi_0$, \textit{rng}); nếu $\pi \notin \textit{seen}$ thì thêm $\pi$ và dừng\;
      nếu hết $A$ lần: thêm \textit{base}(\textit{rng}) với nhãn \texttt{random\_fallback}\;
    }
  }
}
\For{$j \gets 1$ \KwTo $q_{rd}$}{
  $\pi \gets$ \textit{base}(\textit{rng}); sinh lại tối đa $A$ lần nếu trùng; thêm $\pi$ với nhãn \texttt{random}\;
}
\Return{$P$}
\end{algorithm}

Thứ tự gọi bộ sinh ngẫu nhiên được cố định: toàn bộ nhóm \code{nn}, rồi nhóm \code{tw}, rồi nhóm
\code{random}; quần thể không bị xáo lại ở cuối. Bản gốc của mỗi nhóm không tiêu thụ số ngẫu nhiên nào.
Mỗi biến thể tiêu thụ đúng một lần \code{rng.integers} (chọn $k$) và $k$ lần \code{rng.choice} (chọn hai
vị trí). Nhóm \code{random} gọi đúng hàm sinh cá thể của thuật toán gốc,
\code{rng.permutation(np.arange(1, n+1))}. Hệ quả kiểm chứng được bằng test: nếu đặt tỉ lệ
\code{random} = 1, quần thể và trạng thái bộ sinh sau khởi tạo trùng khớp với thuật toán gốc.

\paragraph{Độ phức tạp.} Dựng $\pi_{NN}$ và $\pi_{TW}$ tốn $O(n^2)$, chỉ một lần mỗi lần chạy. Mỗi cá
thể còn lại tốn $O(n)$ để sinh và $O(n)$ để băm khi khử trùng, nhân với số lần thử. Tổng cộng
$O(n^2 + N_{pop}\, n\, A)$ với $A$ là số lần thử tối đa, không đáng kể so với một thế hệ NSGA-II.

\section{Giao diện lập trình}

Mô-đun nằm ở \code{src/improvements/init\_smart.py}:
\begin{Verbatim}
make_initial_population(inst, n_pop, rng, params, baseline_factory=None)
    -> (list[list[int]], list[str])     # hoán vị và nhãn nguồn
make_initial_population_detailed(...)  -> InitResult
    # thêm quota, effective_counts, n_duplicates, n_dedup_retries, notes
\end{Verbatim}
\code{nsga2.py} gọi qua lớp dán mỏng \code{smart\_initial\_population} trong
\code{src/improvements/tv2\_hooks.py}, chỉ khi \code{improvements.init\_smart} bật.
